\section{The differences that matter}
\label{sec:differences}

Section~\ref{sec:dictionary} showed how much is shared. This section isolates
what is not, and argues that the divergences are not accidents of
implementation but a coherent set of trade-offs, each one bought with a
specific currency.

There are six. The first four are deep; the last two are consequences.

\subsection{Difference 1: one sort, not two}
\label{sec:one-sorted}

\fca{} is built on a strict separation. $G$ and $M$ are different sets; an
object is never an attribute; the whole apparatus of primes depends on knowing
which side of the table you are on. \odag{} has one sort: nodes, all the way
down.

The separation buys \fca{} something real --- the Galois connection is a
statement about \emph{two} sets, and the duality it creates is elegant and
generative (every theorem comes with a mirror image). What it costs is
awkwardness in exactly the place ontologies live.

\paragraph{Where it hurts: attributes have no hierarchy.} In
Table~\ref{tab:context} there is no way to say that \textsf{Flight} is a kind
of \textsf{Travel}. \textsf{Flight} is an attribute; \textsf{Travel} would
have to be another attribute; and nothing in the formalism relates two
attributes. The standard workarounds all move the problem rather than solve
it:

\begin{itemize}[leftmargin=1.4em]
\item add \textsf{Travel} as a redundant column and mark it wherever
  \textsf{Flight} or \textsf{Hotel} is marked --- the hierarchy is now encoded
  but not \emph{stated}, and nothing prevents an inconsistent row;
\item use conceptual scaling, which is really the same trick with a name and a
  discipline;
\item move to a many-sorted extension such as Relational Concept Analysis
  \citep{rouanehacene2013}, which is powerful but is a substantially larger
  formalism with an iterative fixpoint computation.
\end{itemize}

\odag{} gets attribute hierarchies for nothing, because \textsf{Flight} is a
node and \cmd{put Flight Travel} is the same operation as
\cmd{put outbound.pdf Flight}. The corollary --- and it is a real consequence,
not a slogan --- is that \emph{the same query primitive traverses both
levels}: \cmd{get(Travel)} returns \textsf{Flight}, \textsf{Hotel} and all the
documents beneath them, without a join, a second index, or a special rule.

\paragraph{What it costs \odag{}.} No duality. There is no ``dual query'' in
\odag{} corresponding to $A'$ --- ``the categories all these items share'' is
computable but is not a first-class operation with its own optimiser, and no
theorem comes for free in mirror image. More importantly, \odag{} cannot
distinguish, in its data model, between ``this thing is an individual'' and
``this thing is a class''. Whether \texttt{outbound.pdf} may acquire children
is a social question, not a checkable one. \fca{} users get that check free.

\begin{keyidea}
The one-sorted design is \odag{}'s largest departure from \fca{} and its most
consequential. It buys hierarchies of categories with no extra machinery, at
the cost of the object/attribute duality and of any built-in class/individual
discipline.
\end{keyidea}

\subsection{Difference 2: asserted, not closed}
\label{sec:not-closed}

\fca{} materialises every concept the data supports. \odag{} materialises
none beyond what was said. Theorem~\ref{thm:dm} says the difference is exactly
the Dedekind--MacNeille completion; Section~\ref{sec:explosion} says what the
completion costs.

The argument for closing is strong and should be stated fairly. Closure
\emph{finds things}. The four unnamed meets of the running example are real
categories that nobody thought to write down, and in a scientific dataset
those are often the point --- that is the entire premise of
\citet{endres2008}'s use of \fca{} on neural data.

The arguments against closing, for a system with \odag{}'s job description,
are four:

\begin{enumerate}[leftmargin=1.7em]
\item \textbf{Size.} Exponential in the worst case, and the worst case is a
  perfectly ordinary table (\S\ref{sec:explosion}). A store that must fit in a
  browser tab and be hashed after every edit cannot have a representation
  whose size is $\#\mathrm{P}$-hard to predict.
\item \textbf{Coincidence.} \fca{} creates a concept for any co-occurrence,
  including accidental ones. With six documents, ``booked flights to Japan''
  is meaningful; with six thousand, most closed sets are noise. \fca{} has no
  notion of a concept not being worth having --- which is precisely the
  observation that motivates iceberg lattices, stability, and the
  MDL programme of Section~\ref{sec:learn-mdl}.
\item \textbf{Instability under edit.} Adding one object can create, destroy
  and rearrange many concepts. For a batch analysis this is fine. For a store
  that is edited continuously and hashed after every edit, it means the
  representation churns globally on local changes --- which destroys
  structural sharing, and with it the cheap snapshots of
  Section~\ref{sec:persistence}.
\item \textbf{Authorship.} A derived concept has no author. In a multi-writer
  setting, ``who said this?'' must have an answer (Section~\ref{sec:agents});
  a concept that exists because of a coincidence between two people's
  unrelated contributions has no honest answer to that question.
\end{enumerate}

\begin{keyidea}
\odag{}'s position: the closure is real and often valuable, but it is
\emph{derived} --- so compute it on demand, locally, per query, and never
store it in the shared, hashed, merged graph. This is the same discipline
\odag{} applies to every derived structure it has (descendant counts, cone
summaries, indexes): regenerable, local, evictable, outside the canonical
root.
\end{keyidea}

\subsection{Difference 3: two kinds of canonical form}
\label{sec:two-canonical-forms}

Both formalisms have a uniqueness theorem, and comparing them is the most
precise way to see what \odag{} is optimising for.

\begin{center}\small
\begin{tabularx}{\textwidth}{@{}p{3.3cm} X X@{}}
\toprule
& \textbf{\fca{}} & \textbf{\odag{}}\\
\midrule
canonical object & the concept lattice $\Bcl(\ctx)$, unique up to isomorphism
  (Thm~\ref{thm:basic}); and the Duquenne--Guigues base, unique on the nose
  \citep{guigues1986}
  & the transitive reduction, unique on the nose
  (Thm~\ref{thm:unique-reduction})\\
what it quotients away & nothing --- it is a complete description
  & assertion order, redundant edges, spelling of equal values\\
cost to compute & exponential in the worst case; $\#\mathrm{P}$-complete to
  size \citep{kuznetsov2001}
  & incremental and local: one edge plus its ancestor set\\
cost to maintain & recompute (a new object may restructure everything)
  & delta on the affected region\\
usable as an address? & in principle; in practice too big and too unstable
  & yes: hash the serialisation, 32 bytes\\
\bottomrule
\end{tabularx}
\end{center}

Both are canonical forms. Only one is a canonical form you can \emph{use as an
identifier}, and that is the whole difference.

\begin{keyidea}
\textbf{\odag{} chose the cheapest non-trivial canonical form that is still
semantic.} It is semantic --- not merely syntactic --- because it quotients
away things that do not change meaning: the order you typed facts in, links
you typed that were already implied, and the way you spelled a value
($400\,\mathrm{g}$ versus $2/5\,\mathrm{kg}$). It is cheap because the
quotient is computed by local edge deletion rather than global inference.

Every step up the expressiveness ladder threatens this. Add arbitrary logical
axioms and canonicalising ``up to logical equivalence'' means canonicalising
an entailment closure; the root silently degrades from a fingerprint of
\emph{knowledge} to a fingerprint of \emph{phrasing}. That, rather than
tractability, is the real wall \odag{} is standing behind.
\end{keyidea}

\subsection{Difference 4: concepts with names, and concepts without}
\label{sec:identity}

This is the subtlest difference and, for the applications in Part~IV, the most
important.

In \fca{}, a concept \emph{is} its (extent, intent) pair. It has no other
identity. This has a consequence that is easy to miss: \textbf{concepts do not
survive changes to the data}. Add one document to Table~\ref{tab:context} and
the concept $(\{\texttt{outbound}, \texttt{inbound}\},
\{\textsf{Flight},\textsf{Japan},\textsf{Booked}\})$ may cease to exist ---
not change, cease. If the new document is a third booked Japanese flight, the
pair with that extent is simply not a concept of the new context. Something
very like it exists, with a different extent. Whether it is ``the same
concept'' is a question the formalism cannot answer, because it has no notion
of sameness across contexts.

In \odag{}, a node \emph{is} its name. \textsf{Flight} is \textsf{Flight}
before and after any edit, in your copy and in mine, in this version and in
the version from last year. Names are the identity at every boundary --- the
public API, the serialisation, and anything crossing between two graphs.

Four things follow, and each of them is load-bearing later:

\begin{enumerate}[leftmargin=1.7em]
\item \textbf{Merge is possible at all.} Two graphs can be merged because
  ``\textsf{Flight} here'' and ``\textsf{Flight} there'' are the same thing by
  definition. Two concept lattices over different contexts have no such
  correspondence; there is no canonical way to merge them.
\item \textbf{References are stable.} An external system can hold a reference
  to \textsf{Flight} and it will still mean something next year. A reference
  to a concept-as-a-pair is invalidated by the next edit.
\item \textbf{Provenance is possible.} ``Alice asserted that
  \texttt{outbound.pdf} is under \textsf{Japan}'' is a durable statement about
  a durable subject. ``Alice asserted concept \#7'' is not.
\item \textbf{The names must be agreed.} This is the price. \odag{}'s
  correctness depends on people using the same string for the same idea, and
  it has no mechanism to enforce that. \fca{}, deriving everything from data,
  never has to ask.
\end{enumerate}

\begin{keyidea}
\fca{} derives concepts and therefore cannot name them stably. \odag{} names
concepts and therefore cannot derive them. This is a genuine trade, not an
oversight on either side, and it is the reason the two formalisms end up
serving different masters: \fca{} serves the analyst looking at a fixed
dataset; \odag{} serves the collaborator, the archive and the verifier, all
of whom need to point at the same thing twice.
\end{keyidea}

There is a long historical argument on \odag{}'s side of this trade, which
Section~\ref{sec:keep-names} develops: every attempt to make identifiers carry
their own meaning --- Wilkins's philosophical language, Dewey decimal
classification, Linnaean binomials --- has paid for it in instability when the
classification was revised, and the mature practice in every field that tried
(chemistry is the cleanest case) is to keep \emph{both} an arbitrary stable
key and a derived meaningful code.

\subsection{Difference 5: a fixed table, or a living store}
\label{sec:mergeability}

\fca{} analyses a context. The context is an input; it does not change while
the analysis runs; there is one of it.

\odag{} is a store. It is edited continuously, by several people who may never
communicate, each of whom holds a replica; and it must converge. Almost every
design decision above is downstream of that requirement:

\begin{center}\small
\begin{tabularx}{\textwidth}{@{}p{4.6cm} X@{}}
\toprule
\textbf{Requirement} & \textbf{What it forced}\\
\midrule
edits arrive one at a time & incremental reduction, delta-maintained counts\\
replicas must converge & merge $=$ union $+$ re-reduction: commutative and
  idempotent (Thm~\ref{thm:merge})\\
agreement must be checkable & canonical form $\Rightarrow$ one hash
  (\S\ref{sec:canonical})\\
answers must not shrink & monotone semantics: no negation in the shared store
  (G2)\\
old answers must stay valid & immutable roots; non-monotone questions must
  name one (\S\ref{sec:asof})\\
\bottomrule
\end{tabularx}
\end{center}

Note the third row in particular. In a merging system, \emph{monotonicity} is
not a stylistic preference; it is what makes union a sound merge. A query
whose answer can shrink when new information arrives cannot be answered
consistently by two replicas that have seen different subsets of the news.
This is why \odag{} refuses negation, closed-world readings and aggregation in
the living store --- and it is also why it can offer them safely against a
pinned root, where the world is closed by \emph{scoping} rather than by
decree.

\fca{} has no analogue of any of this, and needs none. But it also means
there is no such thing as ``merging two concept lattices'', and that a
distributed, multi-author \fca{} is not a thing you can build by adding
features to \fca{}. You would build it by doing what \odag{} does --- sharing
the \emph{assertions} and letting each replica close locally.

\subsection{Difference 6: finite crosses, or computed order}

An \fca{} context is a finite table of crosses. \odag{} additionally supports
an order that is \emph{computed from names}: \verb|weight(4kg)| is below
\verb|weight(..5kg)| because $4 \leq 5$, with no edge stored and no node for
the query term (\S\ref{sec:dimensions}).

In \fca{} terms, the closest analogue is conceptual scaling --- turn each
threshold into an attribute. But scaling must choose the thresholds in
advance, which means quantisation error and a fixed granularity, and it makes
the context grow with the number of thresholds. \odag{}'s values are exact
rationals, arbitrary thresholds are free, and the shared structure carries
only the values that are actually used.

The real \fca{} analogue is \emph{pattern structures}
\citep{ganter2001,kaytoue2011}, which \odag{} appears to have reinvented in
one special case. Section~\ref{sec:learn-patterns} takes this up, because it
is the clearest example in the paper of theory that \odag{} could usefully
import.

\subsection{Summary: the two systems side by side}
\label{sec:summary-of-differences}

\begin{table}[h]\centering\small
\caption{The comparison in one table. A reader who skipped Parts~I and~II can
start here.}
\label{tab:summary}
\vspace{1mm}
\begin{tabularx}{\textwidth}{@{}p{3.5cm} X X@{}}
\toprule
& \textbf{Formal Concept Analysis} & \textbf{\odag{}}\\
\midrule
Input & a fixed binary table & a stream of assertions\\
Sorts & two: objects, attributes & one: nodes\\
Output & every concept, as a complete lattice & the asserted order, minimally
  stored\\
Size & up to exponential & linear in what was said\\
Meets & always exist & only if someone asserted one\\
Concept identity & the (extent, intent) pair & a name string\\
Survives an edit? & concepts may cease to exist & names persist\\
Canonical form & the lattice; the stem base & the transitive reduction\\
Canonical form usable as an address? & no (too large, too volatile) & yes:
  32-byte root\\
Merging two of them & undefined & commutative, idempotent, convergent\\
Query & lattice navigation / lookup & one cone intersection, planned\\
Numeric values & conceptual scaling, or pattern structures & exact rationals,
  order computed from names\\
Negation & expressible on the fixed table & refused in the living store;
  allowed against a pinned root\\
Provenance & no place for it & signed records, per claim\\
Verification & re-run the algorithm & Merkle proofs and certificates\\
Learns concepts from data & yes, all of them & no --- deliberately\\
Best at & \emph{discovering} structure in a dataset & \emph{agreeing on}
  structure between parties\\
\bottomrule
\end{tabularx}
\end{table}

The last row is the honest one-line summary. \fca{} is a discovery instrument.
\odag{} is an agreement instrument. Almost every difference above follows from
that, and Part~III is about what each can take from the other without
becoming the other.
